\documentclass[10pt]{extarticle}
\usepackage[a4paper,margin=1.5cm]{geometry}
\usepackage[T1]{fontenc}
\usepackage{amsmath,amssymb}
\usepackage{graphicx}
\usepackage[export]{adjustbox}
\usepackage{multicol}
\usepackage{enumitem}
\usepackage{booktabs,array,tabularx}
\setlength{\columnsep}{0.7cm}
\setlength{\parindent}{0pt}
\setlength{\parskip}{3pt}
\pagestyle{empty}
\makeatletter
% figure registry (figdefs.tex is generated by the extract step)
\newcommand{\qfigdef}[4]{\@namedef{qf@file@#1}{#2}\@namedef{qf@w@#1}{#3}\@namedef{qf@wide@#1}{#4}}
\newcommand{\qf@check}[1]{\@ifundefined{qf@file@#1}{\PackageError{score2}{Unknown figure id #1}{}}{}}
\newcommand{\qf@img}[4][]{\includegraphics[width=#2\dimexpr\@nameuse{qf@w@#4}pt\relax,max width=#3,max height=0.45\textheight,#1]{\@nameuse{qf@file@#4}}}
% Full-width questions: finalize wraps every question that contains a wide figure in \qfullbegin...\qfullend,
% so the WHOLE question (stem, figures, options) is set across the page in reading order, then the two columns
% resume. The question's head - its number, stem and first figure (\qfullsplit marks the end of the head) -
% is kept on one page, so a number is never left at a page foot with its figure overleaf; the rest may break
% across pages normally, so no big gaps are left. Without a \qfullsplit, the whole question is kept together
% if it is under 45% of a page. Never write these yourself.
\newif\ifqf@full
\newif\ifqf@boxed
\newsavebox{\qf@box}
\newcommand{\qf@open}{\global\qf@fulltrue\par\global\qf@boxedtrue
  \setbox\qf@box\vbox\bgroup\hsize\textwidth\linewidth\textwidth}
\newcommand{\qf@emit}[1]{\par\egroup\global\qf@boxedfalse
  \ifdim\dimexpr\ht\qf@box+\dp\qf@box\relax<#1\textheight\noindent\usebox{\qf@box}\par\else\unvbox\qf@box\fi}
\newcommand{\qfullsplit}{\ifqf@boxed\qf@emit{0.85}\fi}
\newcommand{\qf@close}{\ifqf@boxed\qf@emit{0.45}\fi\global\qf@fullfalse}
\newcommand{\qfullbegin}{\end{multicols}\qf@open}
\newcommand{\qfullend}{\qf@close\begin{multicols}{2}}
\newcommand{\qfullstart}{\qf@open}
\newcommand{\qfulldone}{\qf@close}
% \qfig[scale]{id}: block figure, centred, directly where it is written (size set by the tool).
% Fallback if a wide figure is not inside a full-width question: break out of the columns in place - no floats.
\newcommand{\qfig}[2][1]{\qf@check{#2}%
  \ifqf@full\qf@blk{#1}{#2}\else
  \ifnum\@nameuse{qf@wide@#2}=1 \expandafter\@firstoftwo\else\expandafter\@secondoftwo\fi
  {\end{multicols}\begin{center}\qf@img{#1}{\textwidth}{#2}\end{center}\begin{multicols}{2}}%
  {\qf@blk{#1}{#2}}\fi}
\newcommand{\qf@blk}[2]{\par\smallskip{\centering\qf@img{#1}{\linewidth}{#2}\par}\smallskip}
% \qfigi[scale]{id}: inline figure (inside option lists / tables / minipages); never breaks columns.
\newcommand{\qfigi}[2][1]{\qf@check{#2}\qf@img[valign=c]{#1}{\linewidth}{#2}}
\makeatother
% structure macros - lint reads these to check completeness and figure placement
\newcommand{\papersec}[1]{\par\bigskip{\centering\Large\bfseries PAPER-#1\par}\medskip}
\newcommand{\partsec}[2]{\par\medskip\noindent{\large\bfseries PART-#1: #2}\par\smallskip}
\newcommand{\sect}[1]{\par\smallskip\noindent\textbf{#1}\par}
\newcommand{\qn}[1]{\par\smallskip\noindent\textbf{#1.}~}
\newcommand{\paraq}[2]{\par\smallskip\noindent\textbf{Paragraph for Questions #1 and #2}\par}
\newenvironment{qoptions}{\begin{itemize}[leftmargin=1.8em,itemsep=1pt,topsep=2pt,label={}]}{\end{itemize}}
\newcommand{\opt}[2]{\item[(#1)] #2}
% figure registry (generated by score2_tools.py finalize)
\qfigdef{p1_phy_s1_q02_1}{TNAS - Test 05 (09 Jul 2023)-img/p1_phy_s1_q02_1.png}{96.5}{0}
\qfigdef{p1_phy_s1_q04_1}{TNAS - Test 05 (09 Jul 2023)-img/p1_phy_s1_q04_1.png}{129.3}{0}
\qfigdef{p1_phy_s1_q05_1}{TNAS - Test 05 (09 Jul 2023)-img/p1_phy_s1_q05_1.png}{145.2}{0}
\qfigdef{p1_phy_s1_q15_1}{TNAS - Test 05 (09 Jul 2023)-img/p1_phy_s1_q15_1.png}{89.9}{0}
\qfigdef{p1_phy_s1_q16_1}{TNAS - Test 05 (09 Jul 2023)-img/p1_phy_s1_q16_1.png}{118.7}{0}
\qfigdef{p1_phy_s1_q19_1}{TNAS - Test 05 (09 Jul 2023)-img/p1_phy_s1_q19_1.png}{53.1}{0}
\qfigdef{p1_phy_s1_q20_1}{TNAS - Test 05 (09 Jul 2023)-img/p1_phy_s1_q20_1.png}{55.6}{0}
\qfigdef{p1_phy_s2_q01_1}{TNAS - Test 05 (09 Jul 2023)-img/p1_phy_s2_q01_1.png}{108.2}{0}
\qfigdef{p1_phy_s2_q03_1}{TNAS - Test 05 (09 Jul 2023)-img/p1_phy_s2_q03_1.png}{155.8}{0}
\qfigdef{p1_phy_s2_q06_1}{TNAS - Test 05 (09 Jul 2023)-img/p1_phy_s2_q06_1.png}{137.6}{0}
\qfigdef{p1_phy_s2_q07_1}{TNAS - Test 05 (09 Jul 2023)-img/p1_phy_s2_q07_1.png}{103.2}{0}
\qfigdef{p1_phy_s2_q09_1}{TNAS - Test 05 (09 Jul 2023)-img/p1_phy_s2_q09_1.png}{129.3}{0}
\qfigdef{p1_phy_s2_q10_1}{TNAS - Test 05 (09 Jul 2023)-img/p1_phy_s2_q10_1.png}{126.6}{0}
\qfigdef{p1_chem_s1_q03_1}{TNAS - Test 05 (09 Jul 2023)-img/p1_chem_s1_q03_1.png}{68.9}{0}
\qfigdef{p1_chem_s1_q07_1}{TNAS - Test 05 (09 Jul 2023)-img/p1_chem_s1_q07_1.png}{164.1}{0}
\qfigdef{p1_chem_s1_q07_2}{TNAS - Test 05 (09 Jul 2023)-img/p1_chem_s1_q07_2.png}{165.3}{0}
\qfigdef{p1_chem_s1_q14_1}{TNAS - Test 05 (09 Jul 2023)-img/p1_chem_s1_q14_1.png}{139.1}{0}
\qfigdef{p1_chem_s1_q15_A}{TNAS - Test 05 (09 Jul 2023)-img/p1_chem_s1_q15_A.png}{110.4}{0}
\qfigdef{p1_chem_s1_q15_B}{TNAS - Test 05 (09 Jul 2023)-img/p1_chem_s1_q15_B.png}{95.1}{0}
\qfigdef{p1_chem_s1_q15_C}{TNAS - Test 05 (09 Jul 2023)-img/p1_chem_s1_q15_C.png}{136.2}{0}
\qfigdef{p1_chem_s1_q16_A}{TNAS - Test 05 (09 Jul 2023)-img/p1_chem_s1_q16_A.png}{133.0}{0}
\qfigdef{p1_chem_s1_q16_B}{TNAS - Test 05 (09 Jul 2023)-img/p1_chem_s1_q16_B.png}{77.9}{0}
\qfigdef{p1_chem_s1_q16_C}{TNAS - Test 05 (09 Jul 2023)-img/p1_chem_s1_q16_C.png}{114.1}{0}
\qfigdef{p1_chem_s1_q16_D}{TNAS - Test 05 (09 Jul 2023)-img/p1_chem_s1_q16_D.png}{89.4}{0}
\qfigdef{p1_chem_s1_q18_1}{TNAS - Test 05 (09 Jul 2023)-img/p1_chem_s1_q18_1.png}{95.0}{0}
\qfigdef{p1_chem_s2_q08_1}{TNAS - Test 05 (09 Jul 2023)-img/p1_chem_s2_q08_1.png}{124.6}{0}
\qfigdef{p1_chem_s2_q10_1}{TNAS - Test 05 (09 Jul 2023)-img/p1_chem_s2_q10_1.png}{44.9}{0}
\qfigdef{p2_phy_s1_q02_1}{TNAS - Test 05 (09 Jul 2023)-img/p2_phy_s1_q02_1.png}{158.2}{0}
\qfigdef{p2_phy_s1_q04_1}{TNAS - Test 05 (09 Jul 2023)-img/p2_phy_s1_q04_1.png}{134.9}{0}
\qfigdef{p2_phy_s1_q05_1}{TNAS - Test 05 (09 Jul 2023)-img/p2_phy_s1_q05_1.png}{132.8}{0}
\qfigdef{p2_phy_s1_q07_1}{TNAS - Test 05 (09 Jul 2023)-img/p2_phy_s1_q07_1.png}{227.2}{0}
\qfigdef{p2_phy_s2_para8_10_1}{TNAS - Test 05 (09 Jul 2023)-img/p2_phy_s2_para8_10_1.png}{141.5}{0}
\qfigdef{p2_phy_s2_para11_13_1}{TNAS - Test 05 (09 Jul 2023)-img/p2_phy_s2_para11_13_1.png}{101.1}{0}
\qfigdef{p2_phy_s3_q01_1}{TNAS - Test 05 (09 Jul 2023)-img/p2_phy_s3_q01_1.png}{66.3}{0}
\qfigdef{p2_phy_s3_q02_1}{TNAS - Test 05 (09 Jul 2023)-img/p2_phy_s3_q02_1.png}{234.9}{0}
\qfigdef{p2_phy_s3_q02_2}{TNAS - Test 05 (09 Jul 2023)-img/p2_phy_s3_q02_2.png}{92.5}{0}
\qfigdef{p2_phy_s3_q05_1}{TNAS - Test 05 (09 Jul 2023)-img/p2_phy_s3_q05_1.png}{112.3}{0}
\qfigdef{p2_chem_s1_q02_1}{TNAS - Test 05 (09 Jul 2023)-img/p2_chem_s1_q02_1.png}{89.4}{0}
\qfigdef{p2_chem_s1_q02_2}{TNAS - Test 05 (09 Jul 2023)-img/p2_chem_s1_q02_2.png}{88.0}{0}
\qfigdef{p2_chem_s1_q03_1}{TNAS - Test 05 (09 Jul 2023)-img/p2_chem_s1_q03_1.png}{124.2}{0}
\qfigdef{p2_chem_s1_q06_B}{TNAS - Test 05 (09 Jul 2023)-img/p2_chem_s1_q06_B.png}{88.4}{0}
\qfigdef{p2_chem_s1_q07_A}{TNAS - Test 05 (09 Jul 2023)-img/p2_chem_s1_q07_A.png}{44.9}{0}
\qfigdef{p2_chem_s1_q07_B}{TNAS - Test 05 (09 Jul 2023)-img/p2_chem_s1_q07_B.png}{82.3}{0}
\qfigdef{p2_chem_s1_q07_C}{TNAS - Test 05 (09 Jul 2023)-img/p2_chem_s1_q07_C.png}{87.2}{0}
\qfigdef{p2_chem_s1_q07_D}{TNAS - Test 05 (09 Jul 2023)-img/p2_chem_s1_q07_D.png}{56.1}{0}
\qfigdef{p2_chem_s3_q02_1}{TNAS - Test 05 (09 Jul 2023)-img/p2_chem_s3_q02_1.png}{206.3}{0}
\qfigdef{p2_chem_s3_q03_1}{TNAS - Test 05 (09 Jul 2023)-img/p2_chem_s3_q03_1.png}{302.4}{1}
\qfigdef{p2_chem_s3_q04_1}{TNAS - Test 05 (09 Jul 2023)-img/p2_chem_s3_q04_1.png}{304.6}{1}
\qfigdef{p2_chem_s3_q05_1}{TNAS - Test 05 (09 Jul 2023)-img/p2_chem_s3_q05_1.png}{104.1}{0}
\raggedcolumns
\makeatletter\@beginparpenalty=10000\makeatother
\newcolumntype{L}[1]{>{\raggedright\arraybackslash}p{#1}}
\begin{document}
\begin{center}{\Large\bfseries TNAS --- Test 05 (09 Jul 2023)}\end{center}
\begin{multicols}{2}
\raggedcolumns
\papersec{1}
\partsec{1}{PHYSICS}
\sect{SECTION-I}
\qn{1} Which of the following is not in accordance with kinetic theory of gases? (Here $M$ is molecular weight and $K$ is Boltzmann's constant.)
\begin{qoptions}
\opt{A}{A molecule can only have speed $\sqrt{\dfrac{3RT}{M}}$.}
\opt{B}{A molecule can have speed $\sqrt{\dfrac{12RT}{M}}$.}
\opt{C}{Energy associated with a monoatomic molecule is $\dfrac32KT$.}
\opt{D}{Energy associated with a monoatomic molecule is $\dfrac52KT$.}
\end{qoptions}
\qn{2} Three bulbs are connected to a battery as shown if switch is closed. Then the maximum value of $V$ for which none of the bulb get fuse?
\qfig{p1_phy_s1_q02_1}
\begin{qoptions}
\opt{A}{$V=\dfrac{500}{3}$volt}
\opt{B}{$V=\dfrac{1000}{3}$volt}
\opt{C}{$V=\dfrac{200}{3}$volt}
\opt{D}{$V=400$ volt}
\end{qoptions}
\qn{3} A spherical volume has charge density $\rho(r)=\dfrac{\rho r}{R}$ from $R$ to $2R$, where $r$ is distance from centre. There is no charge from $r=0$ to $r=R$ and $r=2R$ to infinity. Then work done to bring a charge $q$ slowly from $r=2R$ to $r=R$ is :-
\begin{qoptions}
\opt{A}{$\dfrac{11\rho R^2q}{24\varepsilon_0}$}
\opt{B}{$\dfrac{13\rho R^2q}{24\varepsilon_0}$}
\opt{C}{$\dfrac{13\rho R^2q}{12\varepsilon_0}$}
\opt{D}{$\dfrac{11\rho R^2q}{12\varepsilon_0}$}
\end{qoptions}
\qn{4} An infinitely long straight wire carrying current $I$, one side opened rectangular loop and a conductor $C$ with a sliding connector are located in the same plane, as shown in the figure. The connector has length $\ell$ and resistance $R$. It slides to the right with a velocity $v$. The resistance of the conductor and the self inductance of the loop are negligible. Then which of the following is correct if $i$ is current in connecting rod and $F$ is force required to move connecting rod with constant velocity $v$.
\qfig{p1_phy_s1_q04_1}
\begin{qoptions}
\opt{A}{$i\propto r$}
\opt{B}{$i\propto\dfrac1r$}
\opt{C}{$F\propto r$}
\opt{D}{$F\propto\dfrac1{r^2}$}
\end{qoptions}
\qn{5} An infinite long wire having linear charge density $\lambda$ is placed parallel to $x$-axis as shown. Find flux passing through the triangular plate placed in $x$-$z$ plane?
\qfig{p1_phy_s1_q05_1}
\begin{qoptions}
\opt{A}{$\dfrac{\lambda a}{8\varepsilon_0}$}
\opt{B}{$\dfrac{\lambda a}{12\varepsilon_0}$}
\opt{C}{$\dfrac{\lambda a}{16\varepsilon_0}$}
\opt{D}{$\dfrac{\lambda a}{4\varepsilon_0}$}
\end{qoptions}
\qn{6} A straight solenoid of length $\ell$ having a single layer of winding of copper wire whose total mass is equal to $m$. The cross-sectional diameter of the solenoid is assumed to be very less than its length. Take resistivity of copper wire to be $\rho$ and density equal to $d$. Then find the ratio of $\dfrac LR$.\\
($L$ = inductance and $R$ = Resistance of solenoid)
\begin{qoptions}
\opt{A}{$\dfrac{\mu_0m}{4\pi\rho d\ell}$}
\opt{B}{$\dfrac{\mu_0m}{3\pi\rho d\ell}$}
\opt{C}{$\dfrac{\mu_0m}{2\pi\rho d\ell}$}
\opt{D}{$\dfrac{\mu_0m}{\pi\rho d\ell}$}
\end{qoptions}
\qn{7} Speed of sound in hydrogen gas is $v_0$ at temperature $T_0$ then the speed of sound in the mixture of helium and oxygen having same number of moles at temperature $T_0$ is:
\begin{qoptions}
\opt{A}{$\dfrac{V_0}{3}$}
\opt{B}{$\sqrt{\dfrac3{14}}V_0$}
\opt{C}{$\sqrt{\dfrac5{42}}V_0$}
\opt{D}{$\sqrt{\dfrac5{21}}V_0$}
\end{qoptions}
\qn{8} One mole of an ideal gas with adiabatic exponent $\gamma=1.5$ goes through a polytropic process at the end of which the absolute temperature of gas increases 2 times. The polytropic constant of gas equal 2 then increase in entropy of gas in this process. (In SI units)
\begin{qoptions}
\opt{A}{4.76}
\opt{B}{5.76}
\opt{C}{4.26}
\opt{D}{5.26}
\end{qoptions}
\qn{9} Equation of a standing wave is given as $y=2\sin\left(2\pi x+\dfrac\pi6\right)\sin\left(200\pi t+\dfrac\pi4\right)$cm, where $x$ is in cm and $t$ is in sec. Find distance between two consecutive nodes in mm.
\begin{qoptions}
\opt{A}{2}
\opt{B}{5}
\opt{C}{4}
\opt{D}{2.5}
\end{qoptions}
\qn{10} A galvanometer whose resistance is $50\Omega$ has 25 divisions on it. When a current of $4\times10^{-4}$ A passes through it, its needle deflects by one division. What resistance in $\Omega$ should be connected to galvanometer to make it a voltmeter of range 2.5 V.
\begin{qoptions}
\opt{A}{2}
\opt{B}{20}
\opt{C}{200}
\opt{D}{0.2}
\end{qoptions}
\qn{11} Two inductors having inductance $L$ and $2L$ are combine in a circuit. Then which of the following is incorrect if mutual inductance of two inductors is $M=L$.
\begin{qoptions}
\opt{A}{$L_{eq}=L$ when mutual inductance is not opposing and connected in parallel.}
\opt{B}{$L_{eq}=L$ when mutual inductance is opposing and connected in parallel.}
\opt{C}{$L_{eq}=L$ when mutual inductance is opposing and connected in series.}
\opt{D}{$L_{eq}=5L$ when mutual inductance is not opposing and connected in series.}
\end{qoptions}
\qn{12} The magnetic needle performs 20 oscillations per minute in a horizontal plane. If the angle of dip is $30^\circ$, then how many oscillations per minute will this needle perform in vertical north-south plane.
\begin{qoptions}
\opt{A}{$20\left[\dfrac23\right]^{1/4}$}
\opt{B}{$20\left[\dfrac32\right]^{1/4}$}
\opt{C}{$20\left[\dfrac43\right]^{1/2}$}
\opt{D}{$20\left[\dfrac43\right]^{1/4}$}
\end{qoptions}
\qn{13} A rod of length $\ell=100$ cm is fixed at 30 cm from both ends and velocity of transverse wave in rod is $V$ cm/s. Then frequency of third overtone is :-
\begin{qoptions}
\opt{A}{$\dfrac{7V}{10}$}
\opt{B}{$\dfrac{7V}{40}$}
\opt{C}{$\dfrac{10V}{70}$}
\opt{D}{$\dfrac{40V}{70}$}
\end{qoptions}
\qn{14} One end of thin aluminium rod (cross-sectional Area $=10^{-6}$ m$^2$) is welded to the end of copper rod with same diameter. Both rods are very long and under uniform tension of 900 N.\\
(Given that: $Y_{(Cu)}=1.6\times10^{11}$ Nm$^{-2}$; $Y(Al)=0.9\times10^{11}$ Nm$^{-2}$; $\rho_{(Cu)}=8.1\times10^3$ Kgm$^{-3}$; $\rho_{(Al)}=2.5\times10^3$ Kgm$^{-3}$)
\begin{qoptions}
\opt{A}{The value of reflection factor $\dfrac{A_r}{A_i}$ for transverse pulse approaching the junction along copper rod is $\dfrac17$.}
\opt{B}{The value of reflection factor $\dfrac{A_r}{A_i}$ for transverse pulse approaching the junction along copper rod is approximately $\dfrac27$.}
\opt{C}{The value of reflection factor $\dfrac{A_r}{A_i}$ for transverse pulse approaching the junction along copper rod is approximately $\dfrac37$.}
\opt{D}{The value of reflection factor $\dfrac{A_r}{A_i}$ for transverse pulse approaching the junction along copper rod is approximately $\dfrac57$.}
\end{qoptions}
\qn{15} Three alternating voltages $v_1=3\sin\omega t$, $v_2=5\sin(\omega t+30^\circ)$ and $v_3=5\sin(\omega t-127^\circ)$ volt connected across a resistance $R=\sqrt{\dfrac73}\,\Omega$ as shown in figure. The value of peak current through resistor is :-
\qfig{p1_phy_s1_q15_1}
\begin{qoptions}
\opt{A}{1A}
\opt{B}{2A}
\opt{C}{3A}
\opt{D}{4A}
\end{qoptions}
\qn{16} For given circuit key is closed. At $t=0$ then which of the following is incorrect?
\qfig{p1_phy_s1_q16_1}
\begin{qoptions}
\opt{A}{In steady state, current through inductor is 0.5A.}
\opt{B}{In steady state potential drop across capacitor is 6V.}
\opt{C}{In steady state energy stored in inductor is 2J.}
\opt{D}{In steady state energy stored in capacitor is 0.36 mJ.}
\end{qoptions}
\qn{17} Which of the following is incorrect regarding induced electric field (symbols have their usual meaning)
\begin{qoptions}
\opt{A}{Work done in moving a test charge in an induced electric field can be zero.}
\opt{B}{Induced electric field is non-conservative.}
\opt{C}{Induced electric field form closed loops.}
\opt{D}{Induced emf in the loop $\varepsilon=\oint\vec E.d\vec\ell=\dfrac{d\phi}{dt}$.}
\end{qoptions}
\qn{18} The equation of transverse wave is given by $\psi=10^{-2}\sin[\pi(30t-x\sqrt3-y)]$ $x,y,\varepsilon,\psi$ are in m; $\varepsilon$ $t$ in sec. If phase difference between two points $A(2\sqrt3,2)$m and $B(3\sqrt3,3)$m is $n\pi$. Find $n$.
\begin{qoptions}
\opt{A}{$\pi$}
\opt{B}{$2\pi$}
\opt{C}{$4\pi$}
\opt{D}{$6\pi$}
\end{qoptions}
\qn{19} Consider two infinite large parallel current carrying sheets. Current per unit width in both sheet is $\dfrac{1}{\sqrt\pi}$A/m. If the direction of current in both sheets is same, then force per unit area on each sheet is :-
\qfig{p1_phy_s1_q19_1}
\begin{qoptions}
\opt{A}{$2\times10^{-7}$ N/m$^2$}
\opt{B}{$2\pi\times10^{-7}$ N/m$^2$}
\opt{C}{$\pi\times10^{-7}$ N/m$^2$}
\opt{D}{$4\times10^{-7}$ N/m$^2$}
\end{qoptions}
\qn{20} A current carrying string is in the shape of ring of radius 1 m. It is placed in a uniform magnetic field perpendicular to plane of the ring. The current is 1A and area of cross section of ring is 0.2 cm$^2$. The density of string is $2\times10^3$ kg/m$^3$. The speed of wave in the string (ring) is 10 m/s, then magnetic field (in T) is :-
\qfig{p1_phy_s1_q20_1}
\begin{qoptions}
\opt{A}{8 T}
\opt{B}{80 T}
\opt{C}{40 T}
\opt{D}{4 T}
\end{qoptions}
\sect{SECTION-II}
\qn{1} For what value of resistance $R$ (in ohm). The power consumed in $R$ is maximum?
\qfig{p1_phy_s2_q01_1}
\qn{2} An inductance $L=10$H and resistance $R=4\Omega$ are connected in series with a battery of emf $E=8$V. Find maximum rate at which energy is stored in the magnetic field. (in SI units)
\qn{3} A particle of mass $m$ is thrown towards a stationary planet with speed $v$. So that it grazes the planet as shown in figure. Find $x$ if $v=\sqrt{\dfrac{xGM}{105R}}$.
\qfig{p1_phy_s2_q03_1}
\qn{4} A force damped oscillator of mass $m$ has a displacement varying with time given by $x=A\sin\omega t$. The resistance force is $-bv$. From this information calculate how much work is done against the resistive force during one cycle of oscillation. [Given (in SI unit) : $b=2$, $\omega=\pi A=1$ \& Take : $\pi^2=10$]
\qn{5} A body is projected vertically upwards from the surface of earth with a velocity just sufficient enough to carry it to infinity. The time taken by it to reach height $h=3R$ is \dots\dots S. (In nearest integer) (Take : $g=10$ m/s$^2$, Radius of earth : $R=6400$ km and $\sqrt2=1.41$)
\qn{6} If the equivalent capacitance between $A$ and $B$ (in $\mu$F) is $x$, then find the value of $10x$.
\qfig{p1_phy_s2_q06_1}
\qn{7} A charge particle $q=4$c and mass $m=0.1$ kg enters in a uniform magnetic field $B=10$ tesla into the plane of paper as shown in figure from point $A$ and leaves at point $C$. If side of Hexagon is $\sqrt3$m. Then speed of charge particles is \dots\dots m/s.
\qfig{p1_phy_s2_q07_1}
\qn{8} In a series LCR circuit with resistance $R$. When only capacitor is removed, current lags behind the voltage by $60^\circ$. When only inductor is removed, current leads the voltage by $30^\circ$. The ratio of impedence to resistance $R$ is $\sqrt{\dfrac\alpha\beta}$ then $\alpha+\beta$ is equal to.
\qn{9} In the circuit shown battery is ideal. The time constant of the circuit is $x$ C. Find the value of $x$.
\qfig{p1_phy_s2_q09_1}
\qn{10} Initially switch $S$ is open at $t=0$. If switch is closed then the total work done by battery 2 will be equal to \_\_\_ $\mu$J.
\qfig{p1_phy_s2_q10_1}
\partsec{2}{CHEMISTRY}
\sect{SECTION-I}
\qn{1} Consider 2 mole of an ideal hydrogen gas at 300K and 5 atm to be expanded adiabatically to a final pressure of 2 atm against constant external pressure 1 atm. Calculate final temperature of gas.
\begin{qoptions}
\opt{A}{150 K}
\opt{B}{200 K}
\opt{C}{270 K}
\opt{D}{600 K}
\end{qoptions}
\qn{2} The enthalpies of combustion of S(s) and $\mathrm{H_2(g)}$ are $-300$ and $-290$ kcal/mol respectively.\par
Given :\par
$\mathrm{SO_3(g)+H_2O(\ell)\longrightarrow H_2SO_4(\ell)}$ $\Delta H=-130$ kcal/mol\par
$\mathrm{SO_2(g)+\frac12O_2(g)\longrightarrow SO_3(g)}$ $\Delta H=-100$ kcal/mol\par
The enthalpy of formation of $\mathrm{H_2SO_4(\ell)}$ would be
\begin{qoptions}
\opt{A}{$-230$ kcal/mol}
\opt{B}{$-620$ kcal/mol}
\opt{C}{$-820$ kcal/mol}
\opt{D}{$-560$ kcal/mol}
\end{qoptions}
\qn{3} 8 gram of an helium (He) gas was taken from state A to state B as shown in given figure. Calculate heat involved in process? ($R=\dfrac{25}{3}$J mol K) :
\qfig{p1_chem_s1_q03_1}
\begin{qoptions}
\opt{A}{5 kJ}
\opt{B}{7.5 kJ}
\opt{C}{2.5 kJ}
\opt{D}{12.5 kJ}
\end{qoptions}
\qn{4} At least how many half-lives should elapse for a $1^{\mathrm{st}}$ order reaction A $\to$ products so that the reaction is at least 95\% completed ? ($\log2=0.3$)
\begin{qoptions}
\opt{A}{4}
\opt{B}{5}
\opt{C}{6}
\opt{D}{7}
\end{qoptions}
\qn{5} At $35^\circ$C, the vapour pressure of $\mathrm{CS_2}$ is 512 mm Hg and that of acetone $\mathrm{CH_3COCH_3}$, is 344 mm Hg. A solution of $\mathrm{CS_2}$ and acetone, in which the mole fraction of $\mathrm{CS_2}$ is 0.25 has a total vapour pressure of 600 mm Hg. Which the following statement about solution of acetone and $\mathrm{CS_2}$ is true ?
\begin{qoptions}
\opt{A}{A mixture of 100 ml of acetone and 100 ml of $\mathrm{CS_2}$ has a total volume of 200 ml.}
\opt{B}{When acetone and $\mathrm{CS_2}$ are mixed at $35^\circ$C, heat is released to surroundings.}
\opt{C}{When acetone and $\mathrm{CS_2}$ are mixed at $35^\circ$C, heat is absorbed from surroundings.}
\opt{D}{Raoult's law is obeyed by both $\mathrm{CS_2}$ and acetone for the solution in which the mole fraction of $\mathrm{CS_2}$ is 0.25.}
\end{qoptions}
\qn{6} If $\lambda_0$ and $\lambda$ be the threshold wavelength and wavelength of incident light, the velocity of photoelectron ejected from the metal surface is
\begin{qoptions}
\opt{A}{$\sqrt{\dfrac{2hc}{m}\left(\dfrac{\lambda_0-\lambda}{\lambda\lambda_0}\right)}$}
\opt{B}{$\sqrt{\dfrac{2h}{m}\left(\dfrac1{\lambda_0}-\dfrac1\lambda\right)}$}
\opt{C}{$\sqrt{\dfrac{2h}{m}(\lambda_0-\lambda)}$}
\opt{D}{$\sqrt{\dfrac{2hc}{m}(\lambda_0-\lambda)}$}
\end{qoptions}
\qn{7} Select the most appropriate option about ligand considering below given reactions : (tu = thiourea)\par
\textbf{Reaction-1}
\qfig{p1_chem_s1_q07_1}
\textbf{Reaction-2}
\qfig{p1_chem_s1_q07_2}
\begin{qoptions}
\opt{A}{In reaction-1 product, thiosulphate act as bidentate ligand.}
\opt{B}{In reaction-2 product, thiosulphate act as bidentate ligand.}
\opt{C}{In reaction-1 product, tu act as bidentate ligand.}
\opt{D}{In reaction-2 product, tu act as bidentate ligand.}
\end{qoptions}
\qn{8} Correct statement is :\par
(i) In $13^{\mathrm{th}}$ group stability of lower oxidation state increases down the group\par
(ii) Pentavalent bismuth act as reducing agent\par
(iii) $\mathrm{Pb^{4+}}$ is good oxidising agent because $\mathrm{Pb^{4+}}$ have tendency to gain electron and convert into more stable $\mathrm{Pb^{2+}}$ due to inert pair effect.\par
(iv) $\mathrm{BiF_3}$ exist where as $\mathrm{BiF_5}$ not exist.
\begin{qoptions}
\opt{A}{(i), (iv)}
\opt{B}{(i), (iii)}
\opt{C}{(ii), (iii)}
\opt{D}{(ii), (iv)}
\end{qoptions}
\qn{9} Match the molecules under list (1) with shapes under list (2) select the correct answers from the sets (A), (B),(C) and (D) -\par
{\footnotesize\setlength{\tabcolsep}{3pt}\begin{tabular}{|l|l|l|l|}
\hline
\multicolumn{2}{|c|}{\textbf{List (1) molecule}} & \multicolumn{2}{c|}{\textbf{List (2) shape}}\\ \hline
(P) & $\mathrm{PCl_5}$ & (i) & Linear\\ \hline
(Q) & $\mathrm{XeF_2}$ & (ii) & Triangular planar\\ \hline
(R) & $\mathrm{SO_3}$ & (iii) & Trigonal bipyramidal\\ \hline
(S) & $\mathrm{PH_3}$ & (iv) & Trigonal pyramidal\\ \hline
 & & (v) & Tetrahedral\\ \hline
\end{tabular}}
\begin{qoptions}
\opt{A}{P - (iv), Q - (i), R - (ii), S - (iii)}
\opt{B}{P - (iii), Q - (i), R - (ii), S - (v)}
\opt{C}{P - (iii), Q - (i), R - (v), S - (iv)}
\opt{D}{P - (iii), Q - (i), R - (ii), S - (iv)}
\end{qoptions}
\qn{10} {\footnotesize\setlength{\tabcolsep}{3pt}\begin{tabular}{|l|l|l|p{0.33\linewidth}|}
\hline
\multicolumn{2}{|p{0.3\linewidth}|}{\textbf{Column-I (Complex ion)}} & \multicolumn{2}{p{0.4\linewidth}|}{\textbf{Column-II (Hybridisation, No. of unpaired electrons)}}\\ \hline
(P) & $\mathrm{[Cr(H_2O)_6]^{3+}}$ & (1) & $sp^3$, 3\\ \hline
(Q) & $\mathrm{[Co(C\ell)_4]^{2-}}$ & (2) & $sp^3d^2$, 5\\ \hline
(R) & $\mathrm{[Ni(NH_3)_6]^{2+}}$ & (3) & $d^2sp^3$, 3\\ \hline
(S) & $\mathrm{[MnF_6]^{4-}}$ & (4) & $dsp^3$, 3\\ \hline
 & & (5) & $sp^3d^2$, 2\\ \hline
\end{tabular}}
\begin{qoptions}
\opt{A}{P$\to$3, Q$\to$1, R$\to$5, S$\to$2}
\opt{B}{P$\to$3, Q$\to$5, R$\to$1, S$\to$2}
\opt{C}{P$\to$1, Q$\to$3, R$\to$5, S$\to$2}
\opt{D}{P$\to$3, Q$\to$1, R$\to$2, S$\to$5}
\end{qoptions}
\qn{11} {\footnotesize\setlength{\tabcolsep}{3pt}\begin{tabular}{|l|l|l|l|}
\hline
\multicolumn{2}{|c|}{\textbf{List-I (Complex)}} & \multicolumn{2}{c|}{\textbf{List-II (Property)}}\\ \hline
(P) & $\mathrm{[Ni(CO)_4]}$ & (1) & Paramagnetic\\ \hline
(Q) & $\mathrm{[Cu(NH_3)_4]SO_4}$ & (2) & Diamagnetic\\ \hline
(R) & $\mathrm{K_3[Co(CN)_6]}$ & (3) & Square planar\\ \hline
(S) & $\mathrm{[Mn(H_2O)_6]Cl_2}$ & (4) & High spin\\ \hline
 & & (5) & Low spin\\ \hline
\end{tabular}}\par
Select the most appropriate match :-
\begin{qoptions}
\opt{A}{P$\to$3, Q$\to$1, R$\to$4, S$\to$4}
\opt{B}{P$\to$1, Q$\to$2, R$\to$5, S$\to$5}
\opt{C}{P$\to$2, Q$\to$1, R$\to$5, S$\to$4}
\opt{D}{P$\to$2, Q$\to$3, R$\to$1, S$\to$5}
\end{qoptions}
\qn{12} Select the correct statement(s) about $\mathrm{[Fe(CN)_6]^{3-}}$ :-
\begin{qoptions}
\opt{A}{Total seven linkage isomers possible}
\opt{B}{Optical isomerism not possible in any possible linkage isomer.}
\opt{C}{Total possible isomers = 10}
\opt{D}{All of above are correct}
\end{qoptions}
\qn{13} which of the following overlaps is incorrect (assuming $Z$-axis is internuclear axis)?\par
(i) $2p_y+2p_y\to\pi$-Bond formation\par
(ii) $2p_x+2p_x\to\sigma$-Bond formation\par
(iii) $3d_{xy}+3d_{xy}\to\pi$-Bond formation\par
(iv) $2s+2p_y\to\pi$-Bond formation\par
(v) $3d_{xy}+3d_{xy}\to\delta$-Bond formation\par
(vi) $2p_z+2p_z\to\sigma$-Bond formation
\begin{qoptions}
\opt{A}{(i), (ii), (iii)}
\opt{B}{(iii), (vi)}
\opt{C}{(ii), (v)}
\opt{D}{(ii), (iii), (iv)}
\end{qoptions}
\qn{14} Consider following reaction
\qfig{p1_chem_s1_q14_1}
Choose incorrect statements?
\begin{qoptions}
\opt{A}{Degree of unsaturation of B is three}
\opt{B}{Final product `D' is optically inactive}
\opt{C}{Four hydroxy group is present in `D'}
\opt{D}{`C' have no chiral center.}
\end{qoptions}
\qn{15} Choose reaction with correct major product?
\begin{qoptions}
\opt{A}{\qfigi{p1_chem_s1_q15_A}}
\opt{B}{\qfigi{p1_chem_s1_q15_B}}
\opt{C}{\qfigi{p1_chem_s1_q15_C}}
\opt{D}{$\mathrm{CCl_4\xrightarrow[HF+SbF_5]{2AgF\ or}freon\ 12}$}
\end{qoptions}
\qn{16} In which of the following reaction the product will also be optically active?
\begin{qoptions}
\opt{A}{\qfigi{p1_chem_s1_q16_A}}
\opt{B}{\qfigi{p1_chem_s1_q16_B}}
\opt{C}{\qfigi{p1_chem_s1_q16_C}}
\opt{D}{\qfigi{p1_chem_s1_q16_D}}
\end{qoptions}
\qn{17} Isopentane $\xrightarrow{\mathrm{Cl_2/h\nu}}$ Monochloro product(s)\par
Select incorrect statement for above reaction?
\begin{qoptions}
\opt{A}{Total monochloro product are six.}
\opt{B}{Total chiral monochloro products are four.}
\opt{C}{Total fractions after fractional distillation are four.}
\opt{D}{In photohalogenation reaction, carbocation and carbanion are present are as intermediate.}
\end{qoptions}
\qn{18} \qfig{p1_chem_s1_q18_1}
Choose correct option?
\begin{qoptions}
\opt{A}{P is optically active.}
\opt{B}{Reaction of P with HBr give optically active product.}
\opt{C}{P can give bromine water test.}
\opt{D}{Reaction of P with $\mathrm{HBr/R_2O_2}$ give one chiral centre containing product.}
\end{qoptions}
\qn{19} L$-$glucose $\xrightarrow[\Delta]{\mathrm{Red\ P\ +\ HI}}$ Q\par
Choose the correct option?
\begin{qoptions}
\opt{A}{`Q' contains five hydroxyl groups}
\opt{B}{L-glucose can give positive Tollen's test.}
\opt{C}{Structure of glucose have total 32 stereo isomers.}
\opt{D}{Monochlorination of ``Q'' gives ``20'' monochloro products.}
\end{qoptions}
\qn{20} Consider following compounds and find correct increasing order of reactivity with Grignard reagent.\par
(a) $\mathrm{Cl-\overset{O}{\overset{\|}{C}}-Cl}$\quad (b) $\mathrm{CH_3COCl}$\quad (c) $\mathrm{Me-\overset{O}{\overset{\|}{C}}-O-H}$\par
(d) $\mathrm{H_3C-\overset{O}{\overset{\|}{C}}-O-\overset{O}{\overset{\|}{C}}-CH_3}$\quad (e) $\mathrm{CH_3COO^{\ominus}}$\quad (f) $\mathrm{H_3CCOOMe}$
\begin{qoptions}
\opt{A}{$a>b>d>f>c>e$}
\opt{B}{$c<e<f<d<b<a$}
\opt{C}{$a>b>c>f>d>e$}
\opt{D}{$e<f<d<b<a<c$}
\end{qoptions}
\sect{SECTION-II}
\qn{1} Maximum work done by the gas in the isothermal expansion of argon (Ar) gas at $27^\circ$C from 4 L to 16 L was equal to 4200 cal. What is the amount (in gram) of argon subjected to such an expansion?\\
(Ar = 40 g/mol, $\ln2=0.7$) ($R=2$ cal/mol-K)
\qn{2} The gas phase decomposition of dimethylether follows first order kinetics\par
$\mathrm{CH_3{-}O{-}CH_3(g)\longrightarrow CH_4(g)+H_2(g)+CO(g)}$\par
The reaction is carried out in constant volume container at $500^\circ$C and has a half-life of 14.5 min. Initially only dimethylether is present. Find $Z$ after 29 min.\par
\textbf{($Z$ = vapour density of mixture $\times$ 10)}
\qn{3} For the reaction $\mathrm{2A(g)+B(g)\longrightarrow C(g)+D(g)}$, at 400 K, following data are observed. Find out the rate constant (in atm$^{-2}$ sec$^{-1}$) of disappearance of A(g).\par
{\footnotesize\setlength{\tabcolsep}{3pt}\begin{tabular}{|c|p{0.17\linewidth}|p{0.17\linewidth}|p{0.17\linewidth}|}
\hline
\textbf{S.No.} & \textbf{Initial pressure of A(g)} & \textbf{Initial pressure of B(g)} & \textbf{rate of reaction (atm/sec)}\\ \hline
1. & 0.1 atm & 0.4 atm & $16\times10^{-3}$\\ \hline
2. & 0.05 atm & 0.4 atm & $4\times10^{-3}$\\ \hline
3. & 0.1 atm & 0.8 atm & $32\times10^{-3}$\\ \hline
\end{tabular}}
\qn{4} The time period of the electron in ground state of H-atom is two times the time period of the electron in first excited state of certain hydrogen like atom (atomic number $z$). The ionization energy of electron (in ev) present in third excited level of hydrogen like atom is :\par
[Given that ground state energy of electron in H-atom $=-13.6$ eV]
\qn{5} How many species have at least 1 lone pairs on its central atom.\par
$\mathrm{PCl_4^+}$, $\mathrm{BrF_5}$, $\mathrm{TeF_5^-}$, $\mathrm{I(CN)_2^-}$, $\mathrm{ClF_3}$, $\mathrm{CH_3^+}$, $\mathrm{CO_2}$
\qn{6} Out of the following. How many have correct IUPAC naming :\par
(1) $\mathrm{[Ni(CN)_4]^{2-}}$ - Tetracyanonickel(II) ion\par
(2) $\mathrm{[Pt(py)_4][PtCl_4]}$ - Tetrapyridineplatinum(II) tetrachorideplatinate(II)\par
(3) $\mathrm{[Ni\ (dmg)_2]}$ - Bis(dimethylglyoximato)nickel(II)\par
(4) $\mathrm{K_3[Fe(CN)_5NO]}$ - Potassium pentacyanonitrosylferrate(II)\par
(5) $\mathrm{[Fe(CO)_5]}$ - Pentacyanocaronylferrate(0)\par
(6) $\mathrm{K_2[HgI_4]}$ - Potassiumtetraiodidomercurate(II)\par
(7) $\mathrm{[Pt(NH_3)_4][CuCl_4]}$ - Tetraammineplatinum(IV) tetrachloridocuprate(II)\par
(8) $\mathrm{[Cu(gly)_2]}$ - Diglycinatecopper(II)\par
(9) $\mathrm{K_4[Fe(CN)_6]}$ - Potassium hexacyanidoferrate(II)\par
(10) $\mathrm{[Pt(NH_3)_6]Cl_4}$ - Hexaammineplatinum(IV) chloride.
\qn{7} In how many of the following cases ratio of sigma bond and pi bond is 1 : 1\par
{\footnotesize\setlength{\tabcolsep}{3pt}\begin{tabular}{|l|l|l|l|}
\hline
(i) $\mathrm{C_2(CN)_4}$ & (ii) $\mathrm{(CN)_2}$ & (iii) $\mathrm{SO_3}$ & (iv) $\mathrm{C_3O_2}$\\ \hline
(v) $\mathrm{CO_2}$ & (vi) $\mathrm{S_3O_9}$ & (vii) $\mathrm{P_4O_6}$ & (viii) $\mathrm{P_4O_{10}}$\\ \hline
\end{tabular}}
\qn{8} How many moles of $\mathrm{CO_2}$ are released in following reaction :
\qfig{p1_chem_s2_q08_1}
\qn{9} If aliphatic aldehyde and ketone of molecular mass = 86 u are separately treated with $\mathrm{CH_3MgBr}$ followed by acid hydrolysis, how many of them will give racemic mixture?
\qn{10} Consider following reaction sequence
\qfig{p1_chem_s2_q10_1}
$\xrightarrow[\mathrm{(ii)\ H_2O}]{\mathrm{(i)\ EtMgCl\ (excess)}}$ (B) $\xrightarrow[\Delta]{\mathrm{Conc.H_2SO_4}}$ (C) (Major) $\xrightarrow{\mathrm{O_3/Zn}}$ F + G\par
If G can give Fehling test, alpha hydrogen present in (C) and (G) is respectively $x$, $y$ then find value of $x-y$?
\partsec{3}{MATHEMATICS}
\sect{SECTION-I}
\qn{1} Let $f:[0,1]\to R$ be a derivable function with $f(0)=0$, then $\displaystyle\lim_{n\to\infty}n^2\int_0^{1/n}f(t)\,dt$ is equal to-
\begin{qoptions}
\opt{A}{$\dfrac{f'(0)}{2}$}
\opt{B}{$f'(0)$}
\opt{C}{$2f'(0)$}
\opt{D}{0}
\end{qoptions}
\qn{2} The complete set of values the parameter `$a$' so that the point $P\left(a,\dfrac{1}{1+a^2}\right)$ doesn't lie outside the triangle formed by the lines $15y=x+1$, $78y=118-23x$ and $y+2=0$ is -
\begin{qoptions}
\opt{A}{$(0,5)$}
\opt{B}{$[2,5]$}
\opt{C}{$(1,5)$}
\opt{D}{$[0,2]$}
\end{qoptions}
\qn{3} Sum of all the values of $t$ for which the lines $2x+3y=5$, $t^2x+ty-6=0$, and $3x-2y-1=0$ are concurrent
\begin{qoptions}
\opt{A}{$\dfrac14$}
\opt{B}{3}
\opt{C}{$-1$}
\opt{D}{2}
\end{qoptions}
\qn{4} $\displaystyle\int\dfrac{2\left(x^3-1\right)}{x\left(2x^3+1\right)}dx$ is equal to- (where $C$ is the constant of integration)
\begin{qoptions}
\opt{A}{$\ell n\left(2x+\dfrac1{x^2}\right)+C$}
\opt{B}{$\ell n\left(2x^3+1\right)+C$}
\opt{C}{$\tan^{-1}\dfrac{1}{x^{3/2}+1}+\ell n\left(\dfrac{x-1}{x+1}\right)+C$}
\opt{D}{$\ell n\left(x+\sqrt{\dfrac{x-1}{x+1}}\right)+C$}
\end{qoptions}
\qn{5} If the function $g(x)=\begin{cases}ae^x,&x\leqslant0\\b\cos x+x,&x>0\end{cases}$ is differentiable, then the value of $a^2+b^2$ is -
\begin{qoptions}
\opt{A}{5}
\opt{B}{2}
\opt{C}{1}
\opt{D}{13}
\end{qoptions}
\qn{6} The number of real solutions of the equation $\tan(e^x)=e^x+e^{-x}$, $x>0$, is
\begin{qoptions}
\opt{A}{0}
\opt{B}{1}
\opt{C}{2}
\opt{D}{none of these}
\end{qoptions}
\qn{7} If the normal to the curve $y=f(x)$ at the point $(4,6)$ makes an angle $\dfrac{2\pi}{3}$ with the positive $x$ axis in anticlockwise direction, then $f'(4)$ is
\begin{qoptions}
\opt{A}{$\sqrt3$}
\opt{B}{$-\sqrt3$}
\opt{C}{$\dfrac1{\sqrt3}$}
\opt{D}{$\dfrac{-1}{\sqrt3}$}
\end{qoptions}
\qn{8} $\displaystyle\lim_{x\to0}\dfrac{5\sin x+x\cos x}{2\tan x-x^2}$ is
\begin{qoptions}
\opt{A}{$\dfrac52$}
\opt{B}{3}
\opt{C}{0}
\opt{D}{does not exist}
\end{qoptions}
\qn{9} If linear functions $f(x)$ and $g(x)$ satisfy $\displaystyle\int[(1-2x)\cos x+(3+2x)\sin x]\,dx=f(x)\sin x+g(x)\cos x+C$ (where $C$ is constant of integration), then
\begin{qoptions}
\opt{A}{$f(x)=3-2x$}
\opt{B}{$g(x)=5-2x$}
\opt{C}{$f(x)=3+2x$}
\opt{D}{$g(x)=-5+2x$}
\end{qoptions}
\qn{10} Let $f:R\to R$, $f(x)=ax^3+bx^2+cx+d$, has no extreme value, then which of the following is always correct ?
\begin{qoptions}
\opt{A}{$3a+2b+c\geq0$}
\opt{B}{$c\leq0$}
\opt{C}{$(3a+2b+c)c\geq0$}
\opt{D}{$3a+2b+2c\geq0$}
\end{qoptions}
\qn{11} If $f(x)=3-x^2$, $1\leq x\leq4$, then the domain of $\log_e(f(2x))$ is
\begin{qoptions}
\opt{A}{$\left(\dfrac{-\sqrt3}{2},\dfrac{\sqrt3}{2}\right)$}
\opt{B}{$\left[\dfrac12,2\right]$}
\opt{C}{$\left(0,\dfrac{\sqrt3}{2}\right)$}
\opt{D}{$\left[\dfrac12,\dfrac{\sqrt3}{2}\right)$}
\end{qoptions}
\qn{12} If two adjacent sides of a parallelogram are given by $2x^2+2y^2-5xy=0$ and $5x+2y-1=0$ is one of its diagonal, then area of parallelogram is $\dfrac1k$ sq. units, where $k$ is
\begin{qoptions}
\opt{A}{16}
\opt{B}{72}
\opt{C}{18}
\opt{D}{36}
\end{qoptions}
\qn{13} Number of positive integral values of `$k$' for which the equation\par
$k=\left|x+|2x-1|\right|-\left|x-|2x-1|\right|$\par
has exactly three real solutions, is
\begin{qoptions}
\opt{A}{0}
\opt{B}{2}
\opt{C}{3}
\opt{D}{5}
\end{qoptions}
\qn{14} $\displaystyle\int(6x^2+5x+4)(x^2+x+1)^6\cdot x^{27}dx$ equals (where $C$ is integration constant)
\begin{qoptions}
\opt{A}{$\dfrac{x^{28}(1+x+x^2)^7}{7}+C$}
\opt{B}{$\dfrac{x^{24}(1+x+x^2)^7}{7}+C$}
\opt{C}{$\dfrac{x^{24}(1+x+x^2)^6}{7}+C$}
\opt{D}{$\dfrac{x^{28}(1+x+x^2)^8}{7}+C$}
\end{qoptions}
\qn{15} The value of $\displaystyle\lim_{x\to\infty}x^{\frac13}\left((x+1)^{\frac23}-(x-1)^{\frac23}\right)$ is
\begin{qoptions}
\opt{A}{$\dfrac43$}
\opt{B}{$\dfrac{-1}{3}$}
\opt{C}{$\dfrac13$}
\opt{D}{$\dfrac{-2}{3}$}
\end{qoptions}
\qn{16} In the interval $(0,1)$ maximum value of the function $f(x)=|x\ln x|$ is
\begin{qoptions}
\opt{A}{0}
\opt{B}{1}
\opt{C}{$e$}
\opt{D}{$e^{-1}$}
\end{qoptions}
\qn{17} If the equation $\sin^{-1}(x-1)+\cos^{-1}(x-3)+\tan^{-1}\left(\dfrac{x}{-x^2+2}\right)=m$ holds, then the value of `$m$' is
\begin{qoptions}
\opt{A}{$\dfrac{-3\pi}{4}$}
\opt{B}{$\dfrac{3\pi}{4}$}
\opt{C}{$\dfrac{-5\pi}{4}$}
\opt{D}{$\dfrac{5\pi}{4}$}
\end{qoptions}
\qn{18} The expression $x^3-3x^2-9x+c$ can be written in the form $(x-a)^2(x-b)$ if the values of $c$ is
\begin{qoptions}
\opt{A}{$-5$ or 27}
\opt{B}{5 or $-27$}
\opt{C}{5 or 27}
\opt{D}{$-5$ or $-27$}
\end{qoptions}
\qn{19} If $I_1=\displaystyle\int_0^1\dfrac{e^x}{1+x}dx$ and $I_2=\displaystyle\int_0^1\dfrac{x^2}{e^{x^3}(2-x^3)}dx$, then value of $\dfrac{I_1}{I_2}$ is -
\begin{qoptions}
\opt{A}{$e$}
\opt{B}{$\dfrac{3e}{2}$}
\opt{C}{$\dfrac{5e}{2}$}
\opt{D}{$3e$}
\end{qoptions}
\qn{20} If $A$ and $B$ are two points in first quadrant lying respectively on the lines $y=4x$ and $y=2x$ at a distance 4 from origin, then the point where angle bisector of $\angle AOB$ meets $AB$, is
\begin{qoptions}
\opt{A}{$\left(2\left(\dfrac1{\sqrt{17}}+\dfrac1{\sqrt5}\right),\dfrac4{\sqrt{17}}+\dfrac4{\sqrt5}\right)$}
\opt{B}{$\left(2\left(\dfrac1{\sqrt{17}}+\dfrac1{\sqrt5}\right),\dfrac8{\sqrt{17}}+\dfrac8{\sqrt5}\right)$}
\opt{C}{$\left(2\left(\dfrac1{\sqrt{17}}+\dfrac1{\sqrt5}\right),\dfrac8{\sqrt{17}}-\dfrac4{\sqrt5}\right)$}
\opt{D}{$\left(2\left(\dfrac1{\sqrt{17}}+\dfrac1{\sqrt5}\right),\dfrac8{\sqrt{17}}+\dfrac4{\sqrt5}\right)$}
\end{qoptions}
\sect{SECTION-II}
\qn{1} If $m$ is the least value of $f(x)=\sqrt{x^2-2x+2}+\sqrt{x^2-4x+29}$, then $[m]$ is equal to\\
(where $[.]$ denotes greatest integer function)
\qn{2} If `$a$' and `$b$' are the greatest and least values of $f(x)=18x^{4/3}-9x^{1/3}+1$ in the interval $[-1,8]$ then the value of $(a+112b)$ is
\qn{3} Let {\small$I=\displaystyle\int_{-3}^{-2}\left(\dfrac{x^2-x}{x^3-3x+1}\right)^2\left(1+\dfrac1{x^2}+\dfrac1{(1-x)^2}\right)dx$} then the greatest integer less than the value of $I$, is
\qn{4} Straight line $L_1$ is parallel to the bisector of first and third quadrant, forms a triangle of area 2 square units with coordinate axes in second quadrant. Line $L_2$ passes through $M(1,1)$ and has positive $x$ and $y$ intercepts. $L_2$ makes a triangle of minimum area with coordinates axes. Find the area of triangle formed by $L_1$, $L_2$ and $x$-axis.
\qn{5} If $f:[0,\infty)\to[1,\infty)$ is a quadratic function which is invertible and satisfying the equation $\left(f'(x)-4x\right)^2+\left(f''(x)-\lambda\right)^2=f'''(x)$ $\forall$ $x\geq0$, where $\lambda$ is constant then find the value of $\dfrac{15}{\lambda}\displaystyle\int_1^3xf^{-1}(x)\,dx$.
\qn{6} $f(x)=\begin{cases}(x+1)^2&-3<x\leqslant0\\2x^{1/3}+a&0<x\leqslant1\\2-x&1<x<2\end{cases}$ has maxima at $x=0,1$. If number of integral values of `$a$' is $L$ and total number of extrema of $f(x)$ is $M$ then $(L+M)$ is
\qn{7} If $f(x)$, $g(x)$, $h(x)$ be such that $f'(x)=g(x)$, $g'(x)=h(x)$, $h'(x)=f(x)$, $f(0)=4$, $g(0)=5$, $h(0)=0$ then the value of $\left(f^3(x)+g^3(x)+h^3(x)-3f(x).g(x).h(x)\right)$ at $x=2012$ is
\qn{8} A line $L$ having positive slope passing through the point $(1,1)$ intersects the curves $2x^2-y^2+2x-4y+2=0$ at points $A$ and $B$. If the lines joining origin and the points $A$ and $B$ are such that the co-ordinate axes are the angle bisector between them then the equation of line $L$ is given by $ax+by+1=0$, then value of $(3a-b)$ is
\qn{9} If $(1+ax+bx^2)^4=a_0+a_1x+a_2x^2+\dots+a_8x^8$ where $a,b,a_0,a_1,\dots,a_8\in R$ and $a_0+a_1+a_2\neq0$ and $\begin{vmatrix}a_0&a_1&a_2\\a_1&a_2&a_0\\a_2&a_0&a_1\end{vmatrix}=0$, then value of $5\left(\dfrac ab\right)$.
\qn{10} If $\displaystyle\int\dfrac{\sqrt{e^{7x}}}{\sqrt[4]{e^{2x}+e^{-2x}}}dx=\dfrac1k\left(e^{\alpha x}+1\right)^{\beta/\gamma}+C$ (where $C$ is constant of integration \& $k,\alpha,\beta,\gamma$ are positive integers and $\beta,\gamma$ are relatively prime) then $\left(\dfrac k\beta+\dfrac\gamma\alpha\right)$ is
\papersec{2}
\partsec{1}{PHYSICS}
\sect{SECTION-I(i)}
\qn{1} Given are two tuning forks near one another, one of them is of unknown frequency and the other is of frequency 591 Hz. We can hear beats of maximal intensity $I_0$, 5 times each second. At $t=0$ we hear a maxima. Then :
\begin{qoptions}
\opt{A}{Unknown frequency can be 596 Hz}
\opt{B}{Unknown frequency can be 586 Hz}
\opt{C}{Intensity at time $t=2.7$ sec is $I_0$}
\opt{D}{Intensity at time $t=\dfrac{27}{20}$sec is $\dfrac{I_0}{2}$}
\end{qoptions}
\qn{2} Figure shows an acoustic interferometer used to study interference of sound $S$ is a source of sound and $D$ is the detector of sound. Path $SBD$ can be varied in length but path $SAD$ is fixed. It is found that sound intensity has minimum value of 10 $\mu$W/m$^2$ at one position of $B$ and maximum value of 90 $\mu$W/m$^2$ at a position 1.65 cm from the first. Speed of sound in air is 330 m/s.
\qfig{p2_phy_s1_q02_1}
\begin{qoptions}
\opt{A}{The frequency of sound can be 5000 Hz.}
\opt{B}{The frequency of sound can be 15000 Hz.}
\opt{C}{The ratio of amplitude of waves from $SAD$ and $SBD$ is 2 : 1.}
\opt{D}{The ratio of amplitude of waves from $SAD$ and $SBD$ is 3 : 1.}
\end{qoptions}
\qn{3} A horizontal tank of length $2L$ is divided into two equal parts by a thermally insulating piston. Both parts contain a sample of ideal diatomic gas of $n$ moles each at temperature $T_0$. To each side of piston a ideal spring of spring constant $k$ is attached such that the other ends of spring are connected by the walls of the tank. The springs are in natural length when the piston is at the center of the tank. When the gas in right part is slowly heated, the piston moves a distance $\dfrac L2$ towards the left. During the process, the left part is connected to a heat reservoir at temperature $T_0$. Given $k=\dfrac{nRT_0}{L^2}$.
\begin{qoptions}
\opt{A}{The pressure of the gas in the right part when the piston has moved a distance $\dfrac L4$ from the center is $\dfrac{11}{6}\dfrac{nRT_0}{AL}$.}
\opt{B}{The pressure of the gas in the right part when the piston has moved a distance $\dfrac L2$ from the center is $\dfrac{3nRT_0}{AL}$.}
\opt{C}{The work done by gas in the right part on the piston is $0.85\ nRT_0$.}
\opt{D}{The total heat supplied to the right part of the gas is $9.7\ nRT_0$.}
\end{qoptions}
\qn{4} A circuit is made from a wire having resistance per unit length $r/a$, where $a$ is one side of the square as shown in the diagram. A capacitor $C$ is connected into the circuit. A uniform magnetic field exists in the region and increases at a constant rate such that $\dfrac{dB}{dt}=k$. In steady state,
\qfig{p2_phy_s1_q04_1}
\begin{qoptions}
\opt{A}{The current through wire $AD$ is $\dfrac9{23}\dfrac{ka^2}{r}$.}
\opt{B}{Charge on the capacitor is $\dfrac4{23}Cka^2$.}
\opt{C}{Amount of heat released in the time $t$ is $\dfrac{26}{23}\dfrac{k^2a^4t}{r}$.}
\opt{D}{None of the above}
\end{qoptions}
\qn{5} Two loud speakers, A and B, are mounted at ground level, a distance $2d=50.0$ cm apart, and are connected in parallel to the output of an amplifier, which is fed from a variable frequency oscillator. An observer, also at ground level, stands at a point $R$ that is 1.20 km from the line joining the loud speakers and 1.20 km from the perpendicular bisector of $AB$ at $O$. $OR=D$, Figure.  The oscillator is adjusted so that its frequency, $f$, rises linearly with time, at 10 Hz per sec from $f=0$ at time $t=0$. The sound heard by the observer drops to a minimum for the first time when $t=40\sqrt2$ sec.
\qfig{p2_phy_s1_q05_1}
\begin{qoptions}
\opt{A}{The speed of sound is approximately 367 m/s.}
\opt{B}{The speed of sound is approximately 320 m/s.}
\opt{C}{The sound heard by the observer becomes maximum at approximately $t=108.5$ sec.}
\opt{D}{The sound heard by the observer becomes minimum for second time at approximately $t=150$ sec.}
\end{qoptions}
\qn{6} Consider three objects with equal and constant heat capacity at initial temperatures $T_1=225$ K, $T_2=400$ K and $T_3=375$ K. Suppose we wish to raise the temperature of the third object. To do this, we could run a Carnot engine between the first and second objects, extracting work $W$. This work can then be dissipated as heat to raise the temperature of the third object. Even better. it can be stored and used to run a Carnot engine between the first and third object in reverse, which pumps heat into the third object. Assume that all work produced by running engines can be stored and used without dissipation.
\begin{qoptions}
\opt{A}{The minimum temperature $T_L$ to which the first object can be lowered is 300 K.}
\opt{B}{The minimum  temperature $T_2$ to which the first object can be lowered is 250 K.}
\opt{C}{The maximum temperature $T_H$ to which the third object can be raised is 400 K.}
\opt{D}{The maximum temperature $T_H$ to which the third object can be raised is 450 K.}
\end{qoptions}
\qn{7} The schematic diagram of a large cannon is shown in figure. When the gun is fired, high-pressure gases accelerate the projectile inside the barrel to a very high velocity. The reaction force pushes the gun barrel in the direction opposite that of the projectile. Since it is desirable to bring the gun barrel to rest in the shortest time without oscillation, it is made to translate backward against a critically damped spring-damper system called the recoil mechanism. In a particular case, the gun barrel and the recoil mechanism have a mass of 100 kg with a recoil spring of stiffness 10,000 N/m. The gun recoils 0.4 m upon firing. (T
\qfig{p2_phy_s1_q07_1}
\begin{qoptions}
\opt{A}{The damping constant of the damper is 2000 N-s/m.}
\opt{B}{The initial recoil velocity of the gun is 10.8 m/s.}
\opt{C}{The initial acceleration of the gun is 216 m/s$^2$}
\opt{D}{None of the above}
\end{qoptions}
\sect{SECTION-I(ii)}
\paraq{8}{10}
Electric circuit is given, in which battery of emf $10\ V_0$ (volt) can be joined between any of two vertices of the cube each side of cube is identical \& a series combination of a capacitance `$C_0$' and resistance `$R_0$'.
\qfig{p2_phy_s2_para8_10_1}
{\footnotesize\setlength{\tabcolsep}{2pt}\begin{tabular}{|c|p{0.28\linewidth}|c|p{0.1\linewidth}|c|p{0.1\linewidth}|}
\hline
\multicolumn{2}{|c|}{\textbf{Column-1}} & \multicolumn{2}{c|}{\textbf{Column-2}} & \multicolumn{2}{c|}{\textbf{Column-3}}\\ \hline
(I) & Battery is connected between $O$ \& $A$. & (i) & $R_0=84\ \Omega$ & (P) & $C_0=0$ F\\ \hline
(II) & Battery is connected between $F$ \& $E$. & (ii) & $R_0=0\ \Omega$ & (Q) & $C_0=84$ F\\ \hline
(III) & Battery is connected between $F$ \& $B$. & (iii) & $R_0=12\ \Omega$ & (R) & $C_0=10$ F\\ \hline
(IV) & Battery is connected between $O$ \& $B$. & (iv) & $R_0=20\ \Omega$ & (S) & $C_0=12$ F\\ \hline
\end{tabular}}
\qn{8} In which equivalent resistance of the circuit 49 $\Omega$.
\begin{qoptions}
\opt{A}{(I), (ii), (Q)}
\opt{B}{(I), (i), (P)}
\opt{C}{(III), (iii), (P)}
\opt{D}{(III), (iii), (Q)}
\end{qoptions}
\qn{9} In which equivalent capacitance of circuit 12 F.
\begin{qoptions}
\opt{A}{(III), (ii), (R)}
\opt{B}{(III), (ii), (Q)}
\opt{C}{(II), (i), (R)}
\opt{D}{(I), (ii), (Q)}
\end{qoptions}
\qn{10} In which equivalent resistance is 15 $\Omega$.
\begin{qoptions}
\opt{A}{(IV), (i), (P)}
\opt{B}{(IV), (iv), (S)}
\opt{C}{(II), (iv), (P)}
\opt{D}{(IV), (iii), (R)}
\end{qoptions}
\paraq{11}{13}
Answer Q.11, Q.12 and Q.13 by appropriately matching the information given in the three columns of the following table.\par
In a general YDSE setup, a light of wavelength 600 nm is used. `$d$' represents the distance between two slits, `$D$' represents the distance between screen and slits and `$y$' represents the distance of fringe from central maxima.
\qfig{p2_phy_s2_para11_13_1}
{\footnotesize\setlength{\tabcolsep}{2pt}\begin{tabular}{|c|p{0.1\linewidth}|c|p{0.1\linewidth}|c|p{0.1\linewidth}|}
\hline
\multicolumn{2}{|c|}{\textbf{Column-I $(d)$}} & \multicolumn{2}{c|}{\textbf{Column-II $(D)$}} & \multicolumn{2}{c|}{\textbf{Column-III $(y)$}}\\ \hline
(I) & 1 mm & (i) & 1 m & (P) & 1 mm\\ \hline
(II) & 1.5 mm & (ii) & 1.5 m & (Q) & 1.2 mm\\ \hline
(III) & 2 mm & (iii) & 2 m & (R) & 1.5 mm\\ \hline
(IV) & 2.5 mm & (iv) & 2.5 m & (S) & 2 mm\\ \hline
\end{tabular}}
\qn{11} Which of the following corresponding to the maxima ?
\begin{qoptions}
\opt{A}{(III) (ii) (P)}
\opt{B}{(III) (i) (Q)}
\opt{C}{(IV) (iii) (P)}
\opt{D}{(I) (ii) (S)}
\end{qoptions}
\qn{12} Which of the following corresponding to the minima ?
\begin{qoptions}
\opt{A}{(II) (i) (R)}
\opt{B}{(I) (iii) (Q)}
\opt{C}{(II) (iii) (S)}
\opt{D}{(II) (ii) (S)}
\end{qoptions}
\qn{13} Which of the following corresponds to $5^{\mathrm{th}}$ maxima ?
\begin{qoptions}
\opt{A}{(III) (i) (R)}
\opt{B}{(I) (iii) (R)}
\opt{C}{(II) (i) (P)}
\opt{D}{(II) (iii) (Q)}
\end{qoptions}
\sect{SECTION-III}
\qn{1} All the strings and pulley are ideal (very light as compare to $m$) if transverse disturbance is generated at point $A$ then it's speed will be $3x$ then value of $x$ will be (Let mass per unit length in the strings is $\mu$.
\qfig{p2_phy_s3_q01_1}
\qn{2} In the circuit shown (figure-1), all the switches are open initially and all capacitors are uncharged. Now first only $S_1$ is closed and then opened when steady state is reached. Then $S_2$ is closed and again opened when steady state is reached, after this $S_3$ is closed and opened when again steady state is reached, this process repeated countinuously. After closing sixth switch, the fourth capacitor ($8C$) is removed and connected (say at $t=0$) at an uncharged capacitor of original capacitance $C$ which now has a dielectric (dielectric constant 2) filling half its volume (see figure-2), through a resistor $R$. If final charge on capacitor $B$ when steady state reached in figure-2 is $\dfrac n{19}CV$ . The value of $n$ is
\qfig{p2_phy_s3_q02_1}
\qfig{p2_phy_s3_q02_2}
\qn{3} A monoatomic ideal gas expanded adiabatically to $\eta$ times of its initial volume. The ratio of collisions rate of molecules with unit area of container's wall to the initial rate is $\dfrac{N}{N_0}=\dfrac{1}{\eta^{\frac ab}}$. Find $(a+b)$.
\qn{4} A spring mass system with a spring constant $k=5000$ N/m is subjected to a harmonic force of maximum value of 30 N and frequency 20 Hz. The mass is found to oscillate with an amplitude 2cm. The mass of the system is equal to $\dfrac\alpha{32}$kg. The value of $\alpha$ is. (Neglect damping)
\qn{5} For the circuit shown in the diagram, the potential difference between points $D$ and $E(V_{DE})$ is equal to 7V, between points $D$ and $F(V_{DF})$ is equal to 15 V and between points $E$ and $F$ is equal to 20 V, the value of $V_0$ is equal to $\alpha$, fill the value of $\alpha/5$.
\qfig{p2_phy_s3_q05_1}
\partsec{2}{CHEMISTRY}
\sect{SECTION-I(i)}
\qn{1} For the gaseous reaction $A(g)\to 3B(g)+2C(g)$, following data are given.\par
{\footnotesize\begin{tabular}{|l|c|c|c|}
\hline
Time (minutes) & 10 & 20 & 60\\ \hline
Concentration of $A(g)$ & 10 M & 5 M & 2.5 M\\ \hline
\end{tabular}}\par
Select the correct statement(s).
\begin{qoptions}
\opt{A}{Rate of reaction at 20 minutes is 0.1875 M/sec}
\opt{B}{Rate of appearance of $C$ at 20 minutes is 1.5 M/sec.}
\opt{C}{Order of reaction is 3.}
\opt{D}{Order of reaction is 2.}
\end{qoptions}
\qn{2} For the reaction $A(g)\rightleftharpoons P(g)$, the plot of $[A]$ and $[P]$ with time at temperature $T_1$ and $T_2$ are given below. Assume $\Delta H^\circ$ is independent of temperature. If $T_2>T_1$, the correct statement(s) for the reaction $A(g)\rightleftharpoons P(g)$ is/are:
\qfig{p2_chem_s1_q02_1}
\qfig{p2_chem_s1_q02_2}
\begin{qoptions}
\opt{A}{$q=\Delta H>0$}
\opt{B}{$W=0$ (at constant pressure)}
\opt{C}{$q=\Delta U>0$}
\opt{D}{$q=\Delta H=\Delta U<0$}
\end{qoptions}
\qn{3} Werner synthesized and resolved a polynuclear complex containing no carbon atom, Tris[tetraamine-$\mu$-dihydroxidocobalt(III)]cobalt(III) cation also called Hexol. Select the correct statement(s) regarding given complex Hexol.
\qfig{p2_chem_s1_q03_1}
\begin{qoptions}
\opt{A}{Total 3 octahedral units in one complex ion.}
\opt{B}{Octahedral units joined by common edges between two units.}
\opt{C}{Shows optical isomerism.}
\opt{D}{Total 6 bridging hydroxide ligands.}
\end{qoptions}
\qn{4} Consider the given set of species and select correct statements.\par
$\mathrm{H_2,\ Li_2,\ B_2,\ C_2,\ N_2,\ O_2,\ F_2}$\par
Condition 1 $\Rightarrow$ Diamagnetic molecule\par
Condition 2 $\Rightarrow$ Last $e^-$ present in gerade B.M.O.\par
Condition 3 $\Rightarrow$ On $e^-$ addition stability decreases.
\begin{qoptions}
\opt{A}{Only three species follows all three conditions}
\opt{B}{Only two species follows only two conditions}
\opt{C}{Only one specie follows only one condition}
\opt{D}{Only one specie that do not follow any of above condition}
\end{qoptions}
\qn{5} Select correct option for given options.\par
(i) Glycinate (gly)\quad (ii) Acelylacetanato (acac$^-$)\quad (iii) Dipyridine (dipy)\quad (iv) Butylenediamine (bn)\quad (v) Ethylenediamine (en)\quad (vi) 1,10-phenthroline (o-phen)\quad (vii) Acetate (OAc$^-$)\quad (viii) EDTA$^{4-}$ (Ethylenediaminetetraacetate)\quad (ix) Diethylene triamine (dien)
\begin{qoptions}
\opt{A}{Six ligands form 5 member ring with atleast 2 nitrogen atoms and metal centre}
\opt{B}{One ligands form 5 member ring with two oxygen atom and metal centre}
\opt{C}{Seven ligands form 5 member ring with 2 carbon atoms and metal centre}
\opt{D}{Five ligands already contain 6 member ring.}
\end{qoptions}
\qn{6} $\mathrm{H_3C{-}(CH_2)_3{-}OH}\xrightarrow[\mathrm{NEt_3}]{\text{Mesyl chloride}}(P)\xrightarrow{\mathrm{KCN}}(Q)\xrightarrow{\mathrm{H_3O^+}}(R)\xrightarrow{\text{DIBAL--H}}(S)$\par
Select correct statements.
\begin{qoptions}
\opt{A}{$R$ is $\mathrm{H_3C{-}C(=O){-}O{-}H}$}
\opt{B}{$S$ is \qfigi{p2_chem_s1_q06_B}}
\opt{C}{$Q$ will form aldehyde with Sn + HCl}
\opt{D}{When $(R)$ react with $\mathrm{LiAlH_4}$, molecular weight of product is 88.}
\end{qoptions}
\qn{7} Identify reactions give aromatic product.
\begin{qoptions}
\opt{A}{\qfigi{p2_chem_s1_q07_A}\ $\xrightarrow[\Delta]{(\mathrm{NH_4})_2\mathrm{CO_3}}$}
\opt{B}{\qfigi{p2_chem_s1_q07_B}}
\opt{C}{\qfigi{p2_chem_s1_q07_C}}
\opt{D}{\qfigi{p2_chem_s1_q07_D}}
\end{qoptions}
\sect{SECTION-I(ii)}
\paraq{8}{10}
\par
{\footnotesize\setlength{\tabcolsep}{2pt}\begin{tabular}{|c|p{0.37\linewidth}|c|p{0.14\linewidth}|c|p{0.2\linewidth}|}
\hline
\multicolumn{2}{|c|}{\textbf{Column-I}} & \multicolumn{2}{c|}{\textbf{Column-II}} & \multicolumn{2}{c|}{\textbf{Column-III}}\\ \hline
(a) & $[\mathrm{Pt(NH_3)_4Cl_2}]\mathrm{Br_2}$ & (i) & Inner orbital complex & (p) & Total possible isomers = 3\\ \hline
(b) & $[\mathrm{Co(en)(NH_3)_2Cl_2}]\mathrm{Cl}$ & (ii) & Low spin complex & (q) & Total possible isomers = 6\\ \hline
(c) & $\mathrm{K_2[Ni(CN)_4]}$ & (iii) & Outer orbital complex & (r) & Total possible isomers = 8\\ \hline
(d) & $\mathrm{K_4[Ni(SCN)_6]}$ & (iv) & High spin complex & (s) & Total possible isomers = 10\\ \hline
\end{tabular}}\par
\qn{8} Select the CORRECT combination is:
\begin{qoptions}
\opt{A}{(a), (ii), (p)}
\opt{B}{(b), (iv), (q)}
\opt{C}{(c), (i), (s)}
\opt{D}{(d), (iii), (s)}
\end{qoptions}
\qn{9} Select INCORRECT combination is:
\begin{qoptions}
\opt{A}{(a), (ii), (q)}
\opt{B}{(a), (i), (q)}
\opt{C}{(c), (ii), (q)}
\opt{D}{(c), (i), (r)}
\end{qoptions}
\qn{10} Select CORRECT combination in which maximum number of unpaired electrons present in complex:
\begin{qoptions}
\opt{A}{(b), (i), (p)}
\opt{B}{(c), (iii), (q)}
\opt{C}{(d), (iii), (s)}
\opt{D}{(d), (iv), (s)}
\end{qoptions}
\paraq{11}{13}
\par
{\footnotesize\setlength{\tabcolsep}{2pt}\begin{tabular}{|c|p{0.3\linewidth}|c|p{0.1\linewidth}|c|p{0.2\linewidth}|}
\hline
\multicolumn{2}{|c|}{\textbf{Column-I}} & \multicolumn{2}{c|}{\textbf{Column-II}} & \multicolumn{2}{c|}{\textbf{Column-III}}\\ \hline
(I) & Reversible isothermal expansion of an ideal gas & (i) & $\Delta H=0$ & (P) & $w=-nRT\ln\frac{P_1}{P_2}$\\ \hline
(II) & Reversible adiabatic compression of an ideal gas & (ii) & $T_1=T_2$ & (Q) & $PV^\gamma=$ constant\\ \hline
(III) & Reversible adiabatic expansion of an ideal gas & (iii) & $T_1>T_2$ & (R) & $w=\frac{nR(T_2-T_1)}{\gamma-1}$\\ \hline
(IV) & Isothermal free expansion of an ideal gas & (iv) & $T_2>T_1$ & (S) & $PV=$ constant\\ \hline
\end{tabular}}\par
\qn{11}
\begin{qoptions}
\opt{A}{I, i, Q}
\opt{B}{IV, i, S}
\opt{C}{II, iii, Q}
\opt{D}{III, i, Q}
\end{qoptions}
\qn{12}
\begin{qoptions}
\opt{A}{IV, i, Q}
\opt{B}{II, iv, R}
\opt{C}{I, iii, S}
\opt{D}{IV, ii, P}
\end{qoptions}
\qn{13}
\begin{qoptions}
\opt{A}{IV, ii, P}
\opt{B}{II, iii, R}
\opt{C}{III, iii, R}
\opt{D}{I, ii, Q}
\end{qoptions}
\sect{SECTION-III}
\qn{1} Hydrogen atom in ground state is irradiated by UV light of energy $\dfrac{13.6\times48}{49}$ eV. How many wavelengths are possible in the Balmer region?
\qn{2} Ten compounds (a)--(j) are given below.
\qfig{p2_chem_s3_q02_1}
$X$ = Number of compounds giving ($+$) test with $\mathrm{AgNO_3/NH_4OH}$. Find value of $X$?
\qfullbegin
\qn{3} Total number of reaction which form correct major product?
\qfig{p2_chem_s3_q03_1}
\qn{4} Carboxylic acids (i)--(x) are given below.
\qfig{p2_chem_s3_q04_1}
\qfullsplit
(i) The number of given carboxylic acids which exhibit stereoisomerism \& on heating makes cyclic anhydride is $x$;\par
(ii) The number of given carboxylic acids which does not form any cyclic anhydride is $y$.\par
Report your answer as $x+y$?
\qfullend
\qn{5} Identify total number of reagents used for following conversion?
\qfig{p2_chem_s3_q05_1}
(i) $\mathrm{KMnO_4}$ \quad (ii) $\mathrm{K_2Cr_2O_7}$ \quad (iii) $\mathrm{CrO_3.H_2O}$ \quad (iv) $\mathrm{CrO_3.Py.HCl/CCl_4}$ \quad (v) $\mathrm{SeO_2}$ \quad (vi) $\mathrm{MnO_2}$ \quad (vii) Tollen's reagent \quad (viii) Cu/300$^\circ$C
\partsec{3}{MATHEMATICS}
\sect{SECTION-I(i)}
\qn{1} If $\displaystyle\int\sqrt{\sec^2x+3}\,dx=f(x)+\sqrt3\,g(x)+c$ where $c$ is constant of integration and $f(0)=\ell n2$, $g\!\left(\dfrac\pi2\right)=\dfrac\pi3+1$, then
\begin{qoptions}
\opt{A}{$f(x)$ is an odd function}
\opt{B}{$f\!\left(\dfrac\pi4\right)=\ell n\left(1+\sqrt5\right)$}
\opt{C}{$g\!\left(\dfrac{3\pi}2\right)+g\!\left(\dfrac\pi2\right)=2$}
\opt{D}{$g(x)$ is an odd function}
\end{qoptions}
\qn{2} If $I=\displaystyle\int_{-\pi/4}^{\pi/4}\dfrac{e^{\sec x}|\sin x|\sec^2x}{(1+e^{\sin x})}dx$, then
\begin{qoptions}
\opt{A}{$I$ is an irrational number}
\opt{B}{$I=e^{\sqrt2}-e$}
\opt{C}{$I=3.72$}
\opt{D}{$I$ is a rational number}
\end{qoptions}
\qn{3} If $f(x)$ is differentiable at $x=5$ and $f(5)=1$, $f'(5)=\dfrac1{12e}$, then value of $\displaystyle\lim_{x\to5}(x-5)\int_1^{f^2(x)}\dfrac{(x+t)e^t\,dt}{(x-5)^2}$ is
\begin{qoptions}
\opt{A}{0}
\opt{B}{1}
\opt{C}{$\dfrac54$}
\opt{D}{$\dfrac56$}
\end{qoptions}
\qn{4} Let a point `R' be between $P(1,0)$ and $Q(e,1)$ on the graph of $y=\ell nx$ such that the tangent to the graph at point `R' is parallel to the line through points $P$ and $Q$ then which option is/are correct
\begin{qoptions}
\opt{A}{$x$ co-ordinate of the point `R' is $e-1$}
\opt{B}{$x$ co-ordinate of the point `R' is $\dfrac1{e-1}$}
\opt{C}{if the equation of tangent at $R$ is $y=f(x)$ then $\displaystyle\int_1^ef(x)\,dx>\int_1^e\ell nx\,dx$}
\opt{D}{normal at $P$ for $y=\ell nx$ intersect $y=e^x$ at $(0,1)$}
\end{qoptions}
\qn{5} If $f(x)$ be an identity function in $R$ and $g(x)=\displaystyle\sum_{k=1}^{3}\{f(x)-(2016+k)\}^{-1}$, then
\begin{qoptions}
\opt{A}{$g(x)$ is strictly decreasing in $(2018,2019)$}
\opt{B}{$g(x)$ has two distinct real roots}
\opt{C}{Slope of tangent to the curve $g(x)$ at $x=f(2016)$ is $-\dfrac{49}{36}$}
\opt{D}{$\displaystyle\lim_{x\to-\infty}g(x)=0$}
\end{qoptions}
\qn{6} Let $f,g:[-2,3]\to R$ be continuous function which are twice differentiable on the interval $(-2,3)$. Let the values of $f$ and $g$ at the points $-2,0,1$ and $3$ be as given in the following table :\par
{\footnotesize\begin{tabular}{|c|c|c|c|c|}
\hline
 & $x=-2$ & $x=0$ & $x=1$ & $x=3$\\ \hline
$f(x)$ & 11 & $-4$ & 6 & 16\\ \hline
$g(x)$ & 1 & $-2$ & 0 & 2\\ \hline
\end{tabular}}\par
In each of the intervals $(-2,0)$, $(0,1)$ and $(1,3)$ the function $(f-5g)''$ never vanishes. Then the correct statement(s) is (are)
\begin{qoptions}
\opt{A}{$f'(x)-5g'(x)=0$ has exactly three solutions in $(-2,0)\cup(0,1)\cup(1,3)$}
\opt{B}{$f'(x)-5g'(x)=0$ has exactly one solutions in $(0,1)$}
\opt{C}{$f'(x)-5g'(x)=0$ has at least one solution in $(1,3)$}
\opt{D}{$f'(x)-5g'(x)=0$ has at least two solution in $(-2,0)$}
\end{qoptions}
\qn{7} Let $f(\theta)=\dfrac1{\tan^9\theta}\Big((1+\tan\theta)^{10}+(2+\tan\theta)^{10}+\ldots+(20+\tan\theta)^{10}\Big)-20\tan\theta$. If $\displaystyle\lim_{\theta\to\frac\pi2}f(\theta)$ is $\ell$, then
\begin{qoptions}
\opt{A}{$\ell$ is divisible by 100}
\opt{B}{$\ell$ is divisible by 7}
\opt{C}{$\ell$ is divisible by 30}
\opt{D}{$\ell$ is a prime number}
\end{qoptions}
\sect{SECTION-I(ii)}
\paraq{8}{10}
Column 1, 2 and 3 contains function, intervals of monotonic and number of local maxima, local minima, respectively.\par
{\footnotesize\setlength{\tabcolsep}{2pt}\begin{tabular}{|c|p{0.27\linewidth}|c|p{0.24\linewidth}|c|p{0.12\linewidth}|}
\hline
\multicolumn{2}{|c|}{\textbf{Column 1}} & \multicolumn{2}{c|}{\textbf{Column 2}} & \multicolumn{2}{c|}{\textbf{Column 3}}\\ \hline
(I) & $f:(0,\infty)\to R:f(x)=x^2\ell nx$ & (i) & increasing in $\left(e^{-\frac13},\sqrt e\right)$ & (P) & exactly one local maxima\\ \hline
(II) & $g:(0,\infty)\to R:f(x)=e^2\ell nx$ & (ii) & decreasing in $\left(e^{-\frac12},e^{-\frac13}\right)$ & (Q) & exactly one local minima\\ \hline
(III) & $h:(e^{-\pi/4},e^{2\pi})\to R$: $h(x)=\sin(\ell nx)+\cos(\ell nx)$ & (iii) & decreasing in $\left(\dfrac1{\sqrt e},\sqrt e\right)$ & (R) & no local maxima\\ \hline
(IV) & $p:(0,\infty)\to R:\pi(x)=(x^2-3e)e\,\ell nx$ & (iv) & increasing in $\left(\dfrac1{\sqrt e},\sqrt e\right)$ & (S) & no local minima\\ \hline
\end{tabular}}\par
\qn{8} Which of the following options is the only CORRECT combination ?
\begin{qoptions}
\opt{A}{(I) (iv) (P)}
\opt{B}{(I) (ii) (S)}
\opt{C}{(III) (iv) (P)}
\opt{D}{(III) (i) (R)}
\end{qoptions}
\qn{9} Which of the following options is the only CORRECT combination ?
\begin{qoptions}
\opt{A}{(II) (ii) (R)}
\opt{B}{(II) (i) (S)}
\opt{C}{(IV) (i) (Q)}
\opt{D}{(IV) (iv) (P)}
\end{qoptions}
\qn{10} Which of the following options is the only CORRECT combination ?
\begin{qoptions}
\opt{A}{(I) (i) (S)}
\opt{B}{(II) (iv) (P)}
\opt{C}{(III) (i) (S)}
\opt{D}{(IV) (iii) (R)}
\end{qoptions}
\paraq{11}{13}
Let $f(x)$ be a double differentiable function with continuous second derivative defined on the interval $[0,1]$ satisfying the equation $\displaystyle\int_0^x\sqrt{1-(f''(t))^2}\,dt=\int_0^xf'(t)\,dt\ \forall x\in[0,1]$ and $f(0)=f'(0)=0$. Column-1: contains the information about interval of value of $f(x)$; Column-2: contains information about value of $f(x)$; Column-3: contains information about interval of value of definite integral of $f(x)$; Match the following Column(s)\par
{\footnotesize\setlength{\tabcolsep}{2pt}\begin{tabular}{|c|p{0.2\linewidth}|c|p{0.28\linewidth}|c|p{0.22\linewidth}|}
\hline
\multicolumn{2}{|c|}{\textbf{Column 1}} & \multicolumn{2}{c|}{\textbf{Column 2}} & \multicolumn{2}{c|}{\textbf{Column 3}}\\ \hline
(I) & $f\!\left(\frac12\right)<\frac18$ & (i) & $f\!\left(\frac13\right)=\frac1{18}$ & (P) & $\displaystyle\int_0^1f(x)\,dx>1$\\ \hline
(II) & $f\!\left(\frac12\right)>\frac18$ & (ii) & $f\!\left(\frac14\right)=\frac18$ & (Q) & $\displaystyle\int_0^1f(x)\,dx>\frac16$\\ \hline
(III) & $f\!\left(\frac13\right)<\frac1{18}$ & (iii) & $f\!\left(\frac16\right)=2\sin^2\!\left(\frac1{12}\right)$ & (R) & $\displaystyle\int_0^1f(x)\,dx<\frac16$\\ \hline
(IV) & $f\!\left(\frac14\right)<\frac14$ & (iv) & $f\!\left(\frac18\right)=2\sin^2\!\left(\frac1{16}\right)$ & (S) & $\displaystyle\int_0^1f(x)\,dx<\frac14$\\ \hline
\end{tabular}}\par
\qn{11} Which of the following options is the only CORRECT combination ?
\begin{qoptions}
\opt{A}{(I) (ii) (Q)}
\opt{B}{(II) (i) (S)}
\opt{C}{(IV) (ii) (Q)}
\opt{D}{(III) (iii) (R)}
\end{qoptions}
\qn{12} Which of the following options is the only CORRECT combination ?
\begin{qoptions}
\opt{A}{(I) (iv) (Q)}
\opt{B}{(II) (iii) (S)}
\opt{C}{(II) (iv) (P)}
\opt{D}{(IV) (iii) (S)}
\end{qoptions}
\qn{13} Which of the following options is the only CORRECT combination ?
\begin{qoptions}
\opt{A}{(III) (ii) (Q)}
\opt{B}{(IV) (ii) (P)}
\opt{C}{(I) (iii) (R)}
\opt{D}{(III) (ii) (S)}
\end{qoptions}
\sect{SECTION-III}
\qn{1} The number of integral values of $k$ for which the equation $e^x=\dfrac{k}{x-3}$, has exactly two solutions is-
\qn{2} If $\displaystyle\int\dfrac{(3x^2+2x)dx}{x^6+2x^5+x^4+2x^3+2x^2+5}=\dfrac12\tan^{-1}(g(x))+C$, (Where $C$ is constant of integration), then the value of $(g(1)+g(0))$ is (C is constant of integration)
\qn{3} Let $f(x)$ be a twice differentiable function satisfying $xf(x^2)+(x^2+1)f(x)=x$, $\forall\ x\in R$ then the value of $|9f''(1)|$ is
\qn{4} If $f(x+1)=e^{x^2}f(x)$, where $f(x)$ is differentiable function such that $f'\!\left(\dfrac12\right)=5f\!\left(\dfrac12\right)$, then the value of $\dfrac{f'\left(N+\frac12\right)}{f\left(N+\frac12\right)}-N^2$ is ($N=1,2,\ldots\ldots$)
\qn{5} $S(x)=\displaystyle\int\dfrac{dx}{e^x+8e^{-x}+4e^{-3x}}$ \& $R(x)=\displaystyle\int\dfrac{dx}{e^{3x}+8e^x+4e^{-x}}$ \& $M(x)=S(x)-2R(x)$, then $M(\ell n2)-M(\ell n1)$ is equal to
\end{multicols}
\end{document}
